\subsection*{Notes for VI. OPTIMIZATION PROBLEM FORMULATION}
\paragraph{Objective}
Provide the formal mathematical definition: search space, encoding/decoding, constraints, objectives (single and multiobjective), and what objective measurements are used under multi-fidelity.

\paragraph{Keep}
\begin{itemize}
  \item The formal definitions you already drafted (search space, encoding, constraints, metrics), but consolidate them here.
\end{itemize}

\textbf{TODO:}
\begin{itemize}
    \item \paragraph{Rewrite / move}
\end{itemize}
\begin{itemize}
  \begin{itemize}
      \item Move all “single source of truth” definitions (currently scattered in your protocol) into this section and then reference this section elsewhere to avoid duplication.
  \end{itemize}
\end{itemize}

\begin{itemize}
    \item \paragraph{Write next (core elements)}
\end{itemize}
\begin{itemize}
  \begin{itemize}
      \item Define the search space $\Lambda$, the genotype $x \in \Lambda$, the decoding $D(x)$, and the validity constraints $g(x)\le 0$.
      \item Define objectives:
      \begin{itemize}
      \item single-objective $f(x)$ (e.g. validation error at the final rung),
      \item multiobjective $\mathbf{f}(x)=(f_1(x), f_2(x), \dots)$ (e.g., error, compute proxy).
      \end{itemize}
      \item Define a compute proxy precisely (params/FLOPs): input resolution, counting method, whether estimated or measured.
      \item Define multi-fidelity measurement: which rung’s metric is used for selection/model-learning and why.
      \item Define evaluation noise sources and how you handle them (seeds, retraining, aggregation).
  \end{itemize}
\end{itemize}

\paragraph{Peer-review must-haves}
\begin{itemize}
  \item Unambiguous definitions of objectives and proxies (reviewers often challenge the validity of the proxy).
  \item An explicit elite-selection rule under multi-fidelity (to prevent biased learning).
\end{itemize}

\paragraph{Suggested figure/table}
\begin{itemize}
  \item Table: search space specification (ops, nodes, constraints).
  \item Figure: a working example of encoding/serialization for a cell.
\end{itemize}
